
Large language models achieve strong aggregate pass rates on code generation benchmarks, yet a non-trivial fraction of failures persist after round-0 evaluation. On EvalPlus --- a benchmark that augments the base HumanEval and MBPP test suites by approximately $80\times$ --- GPT-4o-mini achieves round-0 pass@1 of 79.3\% on HumanEval+ and 73.5\% on MBPP+, leaving 34 and 100 failures, respectively. These failures are not stylistic or syntactic errors. A systematic analysis using the ruff linter and mypy type checker demonstrates that static analysis fires on only 32.4\% of HumanEval+ failures and 16.0\% of MBPP+ failures. Two-thirds to five-sixths of failures receive no actionable diagnostic signal from the tools that constitute the standard first-pass repair pipeline.

The dominant failure mode is semantic divergence: the generated code is syntactically valid and type-consistent, but algorithmically incorrect. The model has implemented the wrong algorithm or mishandled a class of edge cases. Asking it to try again without additional information --- blind reprompting --- provides minimal repair signal, because the model lacks evidence that its prior solution was wrong or why it was wrong.

Targeted algorithmic repair requires three distinct pieces of information simultaneously. First, the model needs to know what algorithm it was intended to implement: the formal intent captured in the problem docstring. Without this re-anchoring, the model may regenerate the same incorrect solution. Second, the model needs a concrete behavioral gap: a failing test case that specifies an input and the expected output, establishing exactly where the model's behavior diverges from the specification. Third, the model needs a deviation signal: its own actual incorrect output, which enables it to compare its behavior against the expected output and identify the locus of the algorithmic error. Each component addresses a distinct failure of diagnostic information; no single component is sufficient.

This structured triple --- problem docstring's formal intent, failing test input/expected-output pair, and actual model output --- constitutes what the present work calls \textit{specification-aligned repair}. The design is motivated by converging evidence across independent lines of research. Specification grounding improves correct code generation by 38 percentage points, with docstring context identified as the primary driver \citep{Haeri2026Specification}. Removing docstring context from repair feedback causes severe performance degradation \citep{Dai2025FeedbackEval}. Contrastive input/output pairs outperform single failure messages in automated program repair, fixing 143 of 337 bugs against 124 with single failure messages \citep{Kong2024ContrastRepair}. Content --- including actual incorrect output --- rather than mere re-exposure drives LLM repair improvement, with a measured effect of +18 percentage points over bare code re-exposure (p = 0.00042; Iscan, 2026).

Despite this converging literature, no prior work has directly tested specification-aligned repair against blind reprompting on EvalPlus's own semantic failure population, using a controlled three-condition paired design with pre-registered predictions. This gap is consequential: EvalPlus's augmented test suite is substantially stricter than benchmarks used in prior repair work, and the GPT-4o-mini round-0 failure set represents a natural, unselected population of semantic errors rather than a curated bug corpus.

The present paper reports the infrastructure and pre-registration stage of a two-stage research plan. The infrastructure stage establishes and verifies the data foundations required to execute the mechanism experiments without new baseline API calls. The pre-registration stage commits three predictions before any mechanism data are collected, following the pre-registered open science tradition \citep{Nosek2018}. The infrastructure and pre-registration are reported jointly as a complete scholarly contribution to reproducibility and transparency in LLM repair research.

The contributions of this paper are as follows:

\textbf{Data infrastructure verification.} The 134-problem GPT-4o-mini EvalPlus failure set from h-e1 Run 2 is verified as fully recoverable from archive across four conditions: all 134 failure identifiers are confirmed (C1), all 134 stored incorrect solutions are present in \texttt{solutions\textbackslash \_cache.jsonl} (C2), the EvalPlus API is accessible for all 134 task identifiers (C3), and deterministic test selection via \texttt{plus\textbackslash \_input[0]} is confirmed with 780 augmented test cases for a sample task (C4). All four conditions pass; 5/5 pytest integration tests pass in 2.32 seconds. Six tasks with empty \texttt{plus\textbackslash \_fail\textbackslash \_tests} fields define a 128-task working set representing 95.5\% recovery.

\textbf{Pre-registered three-condition comparative design.} Three predictions are registered prior to mechanism experiment execution: (P1) specification-aligned repair achieves significantly higher round-1 pass@1 than blind reprompting (McNemar one-tailed p < 0.05); (P2) specification-aligned repair achieves significantly higher pass@1 than the round-0 baseline with fix rate at least 15\%; (P3, exploratory) the HumanEval+ fix rate exceeds the MBPP+ fix rate under specification-aligned repair.

\textbf{Methodological contribution.} The four-condition existence verification protocol constitutes a reusable template for pre-experiment reproducibility checks in LLM repair research, enabling repair experiments on archived failure sets without new baseline API calls.

The remainder of the paper is organized as follows. Section 2 reviews related work. Section 3 describes the methodology. Section 4 specifies the experimental setup. Section 5 presents the verification results and pre-registered predictions. Section 6 discusses implications and limitations. Section 7 concludes.

